\documentclass[11pt,a4paper]{coverletter}

\senderblock
  {Radheem Bin Razi}
  {Distributed Systems \& Cloud-Native Engineer}
  {Ilmenau, Germany \textbar\ \href{mailto:sheikh.radheem@gmail.com}{sheikh.radheem@gmail.com} \textbar\ \href{https://radheem.github.io/my-cv}{portfolio}}

\begin{document}

%% ===== ENGLISH =====
\recipient{Digitalization Team (Cloud Applications)}{Hiring Team}{}

\opening{Dear Hiring Team,}

\vspace{8pt}\noindent\textbf{1. Why Digitalization Team?}\par\vspace{3pt}

Observability and automation are the quiet engines that let distributed systems scale without manual firefighting. Building that kind of resilient cloud infrastructure is exactly what draws me to the Software Developer / Cloud Engineer role at Digitalization Team. I thrive in interdisciplinary product teams where defining infrastructure and shipping services happen in the same sprint.

\vspace{8pt}\noindent\textbf{2. Why me?}\par\vspace{3pt}

\bullets

  \item Migrated a high-throughput platform from Dapr to native NATS JetStream and gRPC, eliminating framework overhead and streamlining Kubernetes deployments.

  \item Designed a Terraform-aligned deployment pipeline with Skaffold and Cilium, while integrating an MCP-based toolbox that lets LLM agents interact with platform services via declarative, hot-reloadable tool definitions.

  \item Built an event-driven ETL side-channel using Hatchet and sqlc to route streams into PostgreSQL and Dgraph, ensuring reliable data delivery across distributed clusters.

\bulletsend

I am available for full-time work immediately, can relocate anywhere in Germany within two to three weeks, and hold Blue Card eligibility so your team only needs to issue the standard contract.

\closing{Sincerely,}{Radheem Bin Razi}

\clearpage
\selectlanguage{ngerman}

%% ===== DEUTSCH =====
\recipient{Digitalization Team (Cloud Applications)}{Personalabteilung}{}

\opening{Sehr geehrtes Hiring-Team,}

\vspace{8pt}\noindent\textbf{1. Warum Digitalization Team?}\par\vspace{3pt}

Observability und Automatisierung sind die stillen Triebwerke, die es verteilten Systemen ermöglichen, ohne manuelle Notfallmaßnahmen zu skalieren. Der Aufbau einer solchen resilienten Cloud-Infrastruktur ist genau der Grund, warum mich die Rolle als Software Developer / Cloud Engineer bei Digitalization Team anspricht. Ich entfalte meine Stärken in interdisziplinären Produktteams, in denen die Definition der Infrastruktur und das Ausliefern von Services im selben Sprint stattfinden.

\vspace{8pt}\noindent\textbf{2. Warum ich?}\par\vspace{3pt}

\bullets

  \item Ich habe eine hochdurchsatzstarke Plattform von Dapr auf natives NATS JetStream und gRPC migriert, wodurch Framework-Overhead eliminiert und Kubernetes-Deployments optimiert wurden.

  \item Ich habe eine an Terraform angelehnte Deployment-Pipeline mit Skaffold und Cilium konzipiert und eine MCP-basierte Toolbox integriert, die es LLM-Agenten ermöglicht, über deklarative, hot-reloadable Tool-Definitionen mit Plattform-Services zu interagieren.

  \item Ich habe einen ereignisgesteuerten ETL-Nebenkanal mit Hatchet und sqlc aufgebaut, um Datenströme nach PostgreSQL und Dgraph zu routen und so eine zuverlässige Datenübertragung über verteilte Cluster hinweg zu gewährleisten.

\bulletsend

Ich stehe ab sofort für eine Vollzeitbeschäftigung zur Verfügung, kann innerhalb von zwei bis drei Wochen an jeden beliebigen Ort in Deutschland umziehen und erfülle die Voraussetzungen für die Blue Card, sodass Ihr Team lediglich den Standardarbeitsvertrag ausstellen muss.

\closing{Mit freundlichen Grüßen,}{Radheem Bin Razi}

\end{document}
